\section{Conclusion}


Blockchain infrastructure faces a unique quantum threat due to the permanence of
on-chain data, the static nature of validator keys, and the persistent economic value
protected by cryptographic signatures. The harvest-now-decrypt-later attack model
means the threat is active \emph{today}, not in some hypothetical future when quantum
computers mature.

The economic analysis is unambiguous: the cost of hybrid deployment is negligible
relative to the value at risk. The game theory is clear: hybrid deployment is
strictly dominant---it cannot reduce security and may be the difference between
survival and catastrophic loss. The only losing strategy is to wait for certainty
before acting.

%% ========================================================================
%%  BIBLIOGRAPHY
%% ========================================================================
\begin{thebibliography}{20}

\bibitem{shor1997}
P.~W. Shor, ``Polynomial-Time Algorithms for Prime Factorization and Discrete
Logarithms on a Quantum Computer,'' \emph{SIAM Journal on Computing}, vol.~26,
no.~5, pp.~1484--1509, 1997.

\bibitem{grover1996}
L.~K. Grover, ``A Fast Quantum Mechanical Algorithm for Database Search,''
\emph{Proceedings of the 28th Annual ACM Symposium on Theory of Computing},
pp.~212--219, 1996.

\bibitem{roetteler2017quantum}
M. Roetteler, M. Naehrig, K.~M. Svore, and K. Lauter, ``Quantum Resource Estimates
for Computing Elliptic Curve Discrete Logarithms,'' \emph{Advances in Cryptology --
ASIACRYPT 2017}, pp.~241--270, Springer, 2017.

\bibitem{gidney2021factor}
C. Gidney and M. Eker\aa, ``How to Factor 2048 Bit RSA Integers in 8 Hours Using
20 Million Noisy Qubits,'' \emph{Quantum}, vol.~5, p.~433, 2021.

\bibitem{mosca2018}
M.~Mosca, ``Cybersecurity in an Era with Quantum Computers: Will We Be Ready?''
\emph{IEEE Security \& Privacy}, vol.~16, no.~5, pp.~38--41, 2018.

\bibitem{fips203}
NIST, ``Module-Lattice-Based Key-Encapsulation Mechanism Standard (FIPS 203),''
National Institute of Standards and Technology, August 2024.

\bibitem{fips204}
NIST, ``Module-Lattice-Based Digital Signature Standard (FIPS 204),''
National Institute of Standards and Technology, August 2024.

\bibitem{fips205}
NIST, ``Stateless Hash-Based Digital Signature Standard (FIPS 205),''
National Institute of Standards and Technology, August 2024.

\bibitem{boneh2001bls}
D. Boneh, B. Lynn, and H. Shacham, ``Short Signatures from the Weil Pairing,''
\emph{Journal of Cryptology}, vol.~17, no.~4, pp.~297--319, 2004.

\bibitem{webber2022impact}
S.~E. Webber, S. Patil, D. Linke, et al., ``The Impact of Hardware Specifications on
Reaching Quantum Advantage in the Fault Tolerant Regime,'' \emph{AVS Quantum Science},
vol.~4, no.~1, 2022.

\bibitem{campbell2019}
E.~T. Campbell, B.~M. Terhal, and C. Vuillot, ``Roads Towards Fault-Tolerant Universal
Quantum Computation,'' \emph{Nature}, vol.~549, pp.~172--179, 2017.

\bibitem{preskill2018}
J. Preskill, ``Quantum Computing in the NISQ Era and Beyond,'' \emph{Quantum},
vol.~2, p.~79, 2018.

\end{thebibliography}

